\documentclass[12pt]{article}
\usepackage[T1]{fontenc}
\usepackage[utf8]{inputenc}
\usepackage{lmodern}
\usepackage{textcomp}
\usepackage{enumitem}
\usepackage{geometry}
\geometry{margin=1in}
\begin{document}
\begin{center}
Answer Key - Economics - Undergraduate Level Assessment 9 \
\small (Answers are ordered exactly as in the question paper.)
\end{center}

\section*{Multiple Choice}
\begin{enumerate}[label=\arabic*.]
\item \textbf{Answer:} A
\\ \textit{Options:} A) remain constant and taxes rise.; B) fluctuate and taxes rise.; C) fall and taxes remain constant.; D) rise and taxes fall.; E) and taxes rise.
\\ \textit{Note:} During economic expansions, automatic stabilizers such as progressive tax systems lead to increased tax revenues while government expenditures on social programs typically remain stable, resulting in higher taxes and constant spending.
\item \textbf{Answer:} E
\\ \textit{Options:} A) The monopoly would make an economic profit.; B) The monopoly would break even.; C) The monopoly would increase its prices.; D) The deadweight loss in this market would increase.; E) The deadweight loss in this market would decrease.
\\ \textit{Note:} Regulating a monopoly to produce at the allocatively efficient quantity eliminates the inefficiencies associated with underproduction, thereby reducing the deadweight loss in the market.
\item \textbf{Answer:} A
\\ \textit{Options:} A) A price ceiling will have no effect on the quantity of the good supplied.; B) A price ceiling will cause a shift in the supply curve for the good.; C) A price ceiling always results in a surplus of the good.; D) A price ceiling will increase the total revenue of buyers.; E) An effective price ceiling must be at a price above the equilibrium price.
\\ \textit{Note:} A price ceiling, which is a maximum price set by the government, does not alter the quantity supplied because suppliers will still produce the same amount regardless of the price limit imposed, as long as it is above their cost of production.
\item \textbf{Answer:} A
\\ \textit{Options:} A) The United States begins to import rice to make up for a domestic shortage.; B) The United States begins to export rice to make up for a domestic shortage.; C) The United States begins to import rice to eliminate a domestic surplus.; D) The United States begins to export rice to eliminate a domestic surplus.
\\ \textit{Note:} The correct answer is based on the principle that when the world price of rice is lower than the domestic price, U.S. consumers will prefer the cheaper imported rice, leading to an increase in imports to satisfy domestic demand and address the resulting shortfall in local supply.
\item \textbf{Answer:} A
\\ \textit{Options:} A) expansionary monetary policy by lowering the discount rate.; B) expansionary monetary policy by selling Treasury securities.; C) contractionary monetary policy by raising the discount rate.; D) contractionary monetary policy by lowering the discount rate.
\\ \textit{Note:} The correct answer is based on the key concept that lowering the discount rate makes borrowing cheaper for banks, which can increase the money supply and stimulate economic activity, countering the "crowding-out" effect where government spending limits private sector investment.
\end{enumerate}

\section*{True / False}
\begin{enumerate}[label=\arabic*.]
\setcounter{enumi}{5}
\item \textbf{Answer:} True
\\ \textit{Options:} A) True; B) False
\\ \textit{Note:} An autonomous increase in spending is characterized by its independence from current income levels, meaning it can occur due to factors like consumer confidence or government policy changes, rather than being a response to increased earnings.
\item \textbf{Answer:} True
\\ \textit{Options:} A) True; B) False
\\ \textit{Note:} The statement is true because, despite air being essential for life, its abundance in the environment means that its marginal utility is effectively zero for many people, leading to the conclusion that it is not considered a scarce good in economic terms.
\item \textbf{Answer:} False
\\ \textit{Options:} A) True; B) False
\\ \textit{Note:} A decrease in the price of aluminum, which is a key input in bicycle production, would lower production costs, potentially leading to an increase in the supply of bicycles and therefore an increase in demand due to lower prices for consumers.
\item \textbf{Answer:} False
\\ \textit{Options:} A) True; B) False
\\ \textit{Note:} A Durbin Watson statistic close to zero indicates the presence of positive first-order autocorrelation, not that the autocorrelation coefficient is exactly zero.
\item \textbf{Answer:} True
\\ \textit{Options:} A) True; B) False
\\ \textit{Note:} Expectations of future surpluses lead consumers and businesses to anticipate lower prices, causing them to reduce current spending and investment, which shifts the aggregate demand curve to the left.
\end{enumerate}

\section*{Long Form Response}
\begin{enumerate}[label=\arabic*.]
\setcounter{enumi}{10}
\item \textbf{Answer:} John's financial situation can be analyzed by understanding his break-even point and his income level in relation to his marginal propensity to consume (MPC). With a break-even point of $7,000, this is the amount of income he needs to cover his basic expenses. Given that his current income is $3,000, he is short by $4,000 to reach this break-even point. Since his MPC is 3/4, it indicates that he spends 75% of any additional income; however, in this case, he lacks the necessary income to even reach the break-even threshold. Therefore, John will need to borrow $4,000 to cover his expenses and reach financial stability.
\\ \textit{Note:} correct because it accurately determines that John's income of $3,000 is $4,000 short of the \$7,000 needed to cover his expenses, indicating he must borrow this amount to reach his break-even point, while his MPC of 3/4 further emphasizes his need to manage spending against inadequate income.
\item \textbf{Answer:} The unique ownership of the only original signed copy of The Wealth of Nations by Adam Smith illustrates the principles of supply and demand through its exceptional scarcity, which aligns with the concepts of a downward sloping supply curve. As the supply of this historical artifact is fixed at one, its value increases due to heightened demand from collectors, historians, and economists who recognize its significance. This scenario demonstrates that as the quantity of a good decreases, its price tends to rise, reflecting the inverse relationship depicted in the downward sloping supply curve. Consequently, the rarity of the signed copy not only amplifies its market value but also underscores the fundamental economic principle that limited availability can drive up demand and price.
\\ \textit{Note:} correct because it illustrates how the limited supply of a unique item, such as the only original signed copy of The Wealth of Nations, leads to increased demand and higher prices, which aligns with the principles of a downward sloping supply curve in economics.
\end{enumerate}

\vfill
\noindent\textit{End of Answer Key}
\end{document}